\section{Metaverse Companion}

Zoo NFT animals feature stunning fully 3D-rendered presentations designed to mimic the appearance, personalities, and movements of the critically endangered real-world animals they represent. The key capabilities include:

\begin{itemize}[nosep]
\item \textbf{3D Interaction:} Ability to be interacted with in 3 dimensions on any mobile or desktop device by any user, whether they own an NFT or not, on the Zoo website.

\item \textbf{Augmented Reality (AR):} AR capabilities allow any NFT holder to bring Zoo NFT Animals to life in their physical environments. Further functionality for NFT owners will be available after the release of the iOS/Android AR Companion app.

\item \textbf{Virtual Reality (VR):} VR capabilities allow users who own Zoo NFT animals to view and interact with them while immersed in Virtual Reality, both in first-party software offerings and through potential licensing and partnership opportunities.

\item \textbf{Metaverse Integration:} Ability to be adapted and integrated into the Metaverse, including partnerships in development and endless possibilities as the move to the Metaverse accelerates.

\item \textbf{Store of Value:} Unique ability to become literal stores of value through the innovative way Animal NFTs can be ``fed,'' increasing stored value by collateralizing them with not just \$ZOO tokens, but many other popular cryptocurrencies.

\item \textbf{Growth Mechanics:} Ability to grow, hatching from an NFT Egg as a baby, eventually being able to change physical form, stats, APY, and value according to what they are ``fed.''

\item \textbf{AI Updates:} The capability to be updated as new functionality is made possible, with advances related to Artificial Intelligence and Emotional Intelligence currently in development. These updates will allow NFT Animals to interact with their surroundings in AR and the user in all mediums in innovative and exciting ways.
\end{itemize}

\subsection{Tamagotchi-Inspired Design}

Zoo is a game that ties in finance strategies into a metaphor that is a Tamagotchi-like game. Tamagotchi is a handheld digital pet that was first released in Japan in 1996 by Bandai. The virtual pet could be cared for by inputting commands on a keychain-sized device. The original release was popular, with over 80~million units sold worldwide.

\subsection{Seven Key Differentiators}

The vast majority of NFT gaming projects lack 7 key elements that Zoo Labs provides:

\begin{enumerate}[nosep]
\item An extremely clear path to quick and meaningful ROI through a unique approach to staking and making NFTs into stores of value.
\item The ability to appeal to users of younger age groups and across virtually all demographics due to the universal appeal of animals.
\item Support of the Open-Source movement, allowing products and platforms to be built upon and utilized by an endless number of partners and sponsors.
\item Ability to keep demand high and keep NFTs liquid through pooling mechanisms, high staking rewards, credit cards to leverage Zoo assets, and a continuous stream of upgrades.
\item Liquidity sustainability: for every transaction, there is a tax to benefit the mission and sustainable growth of the Zoo ecosystem---4\% goes to the \DAO{} and 4\% goes back to the Liquidity pool.
\item The Zoo \DAO{} keeps the ecosystem decentralized while leveraging holographic consensus, quadratic voting, aged staking, and weighted PoS. Participants maintain real-time control over contributed assets, mechanized and enforced using smart contract code.
\item Artificially intelligent NFTs: Zoo Labs NFTs use AI to understand emotional intelligence, to be a sort of digital pet for a multi-generational audience to play and interact with from their mobile devices.
\end{enumerate}
